
\documentclass{article}
\usepackage{colortbl}
\usepackage{makecell}
\usepackage{multirow}
\usepackage{supertabular}

\begin{document}

\newcounter{utterance}

\centering \large Interaction Transcript for game `matchit\_ascii', experiment `different\_grid', episode 3 with deepseek{-}v4{-}pro.
\vspace{24pt}

{ \footnotesize  \setcounter{utterance}{1}
\setlength{\tabcolsep}{0pt}
\begin{supertabular}{c@{$\;$}|p{.15\linewidth}@{}p{.15\linewidth}p{.15\linewidth}p{.15\linewidth}p{.15\linewidth}p{.15\linewidth}}
    \# & $\;$A & \multicolumn{4}{c}{Game Master} & $\;\:$B\\
    \hline

    \theutterance \stepcounter{utterance}  
    & & \multicolumn{4}{p{0.6\linewidth}}{
        \cellcolor[rgb]{0.9,0.9,0.9}{
            \makecell[{{p{\linewidth}}}]{
                \texttt{\tiny{[P1$\langle$GM]}}
                \texttt{Sudjeluješ u igri pogađanja u suradnji.} \\
\texttt{Cilj je otkriti jesu li ova mreža i druga slika koju samo ja mogu vidjeti iste.} \\
\texttt{Molim te, prvo opiši svoju mrežu. Zatim ću ja dati svoj opis i možemo si međusobno postavljati pitanja o mrežama kako bismo utvrdili jesu li iste. Sada započni svoj kratki opis mreže s "OPIS:" nakon čega slijedi opis. Ne dodaj ništa drugo. Ovo je tvoja mreža:} \\
\texttt{▢ X X X ▢} \\
\texttt{X ▢ X ▢ X} \\
\texttt{X X ▢ X X} \\
\texttt{X X ▢ X X} \\
\texttt{X X ▢ X X} \\
            }
        }
    }
    & \\ \\

    \theutterance \stepcounter{utterance}  
    & \multicolumn{4}{p{0.6\linewidth}}{
        \cellcolor[rgb]{0.9,0.9,0.9}{
            \makecell[{{p{\linewidth}}}]{
                \texttt{\tiny{[P1$\rangle$GM]}}
                \texttt{OPIS: Mreža je 5x5. Gornji red: prazno, X, X, X, prazno. Drugi red: X, prazno, X, prazno, X. Treći, četvrti i peti red: X, X, prazno, X, X.} \\
            }
        }
    }
    & & \\ \\

    \theutterance \stepcounter{utterance}  
    & & & \multicolumn{2}{p{0.3\linewidth}}{
        \cellcolor[rgb]{0.9,0.9,0.9}{
            \makecell[{{p{\linewidth}}}]{
                \texttt{\tiny{[GM$|$GM]}}
                \texttt{continue} \\
            }
        }
    }
    & & \\ \\

    \theutterance \stepcounter{utterance}  
    & & \multicolumn{4}{p{0.6\linewidth}}{
        \cellcolor[rgb]{0.9,0.9,0.9}{
            \makecell[{{p{\linewidth}}}]{
                \texttt{\tiny{[GM$\rangle$P2]}}
                \texttt{Sudjeluješ u igri pogađanja u suradnji.} \\
\texttt{Cilj je otkriti jesu li ova mreža i druga slika koju samo ja mogu vidjeti iste.} \\
\texttt{Molim te, prvo opiši svoju mrežu. Zatim ću ja dati svoj opis i možemo si međusobno postavljati pitanja o mrežama kako bismo utvrdili jesu li iste. Sada započni svoj kratki opis mreže s "OPIS:" nakon čega slijedi opis. Ne dodaj ništa drugo. Ovo je tvoja mreža:} \\
\texttt{▢ ▢ ▢ ▢ X} \\
\texttt{▢ ▢ ▢ X X} \\
\texttt{▢ ▢ X X X} \\
\texttt{▢ X X X X} \\
\texttt{X X X X X} \\
            }
        }
    }
    & \\ \\

    \theutterance \stepcounter{utterance}  
    & & & \multicolumn{4}{p{0.6\linewidth}}{
        \cellcolor[rgb]{0.9,0.9,0.9}{
            \makecell[{{p{\linewidth}}}]{
                \texttt{\tiny{[GM$\langle$P2]}}
                \texttt{OPIS:} \\
\texttt{Mreža je dimenzija 5x5.} \\
\texttt{Gornji lijevi kut je prazan (▢), a donji desni kut je popunjen (X).} \\
\texttt{Popunjeni X{-}ovi tvore oblik obrnutih stepenica koje se spuštaju od gornjeg desnog kuta prema donjem lijevom kutu.} \\
\texttt{Redovi od vrha prema dnu:} \\
\texttt{1. red: 4 prazna, 1 X} \\
\texttt{2. red: 3 prazna, 2 X} \\
\texttt{3. red: 2 prazna, 3 X} \\
\texttt{4. red: 1 prazno, 4 X} \\
\texttt{5. red: 5 X} \\
            }
        }
    }
     \\ \\

    \theutterance \stepcounter{utterance}  
    & & & \multicolumn{2}{p{0.3\linewidth}}{
        \cellcolor[rgb]{0.9,0.9,0.9}{
            \makecell[{{p{\linewidth}}}]{
                \texttt{\tiny{[GM$|$GM]}}
                \texttt{continue} \\
            }
        }
    }
    & & \\ \\

    \theutterance \stepcounter{utterance}  
    & & \multicolumn{4}{p{0.6\linewidth}}{
        \cellcolor[rgb]{0.9,0.9,0.9}{
            \makecell[{{p{\linewidth}}}]{
                \texttt{\tiny{[GM$\rangle$P2]}}
                \texttt{Ovo je mojOPIS: Mreža je 5x5. Gornji red: prazno, X, X, X, prazno. Drugi red: X, prazno, X, prazno, X. Treći, četvrti i peti red: X, X, prazno, X, X.} \\
\texttt{Sada postavi pitanje kako bi saznao nove aspekte moje mreže koji se možda razlikuju od tvoje mreže. Započni s "PITANJE:" i ne dodavaj ništa drugo.} \\
            }
        }
    }
    & \\ \\

    \theutterance \stepcounter{utterance}  
    & & & \multicolumn{4}{p{0.6\linewidth}}{
        \cellcolor[rgb]{0.9,0.9,0.9}{
            \makecell[{{p{\linewidth}}}]{
                \texttt{\tiny{[GM$\langle$P2]}}
                \texttt{PITANJE:} \\
\texttt{U tvojoj mreži, je li donji lijevi kut (pozicija 5,1) prazan (▢) ili popunjen (X)?} \\
            }
        }
    }
     \\ \\

    \theutterance \stepcounter{utterance}  
    & & & \multicolumn{2}{p{0.3\linewidth}}{
        \cellcolor[rgb]{0.9,0.9,0.9}{
            \makecell[{{p{\linewidth}}}]{
                \texttt{\tiny{[GM$|$GM]}}
                \texttt{continue} \\
            }
        }
    }
    & & \\ \\

    \theutterance \stepcounter{utterance}  
    & & \multicolumn{4}{p{0.6\linewidth}}{
        \cellcolor[rgb]{0.9,0.9,0.9}{
            \makecell[{{p{\linewidth}}}]{
                \texttt{\tiny{[P1$\langle$GM]}}
                \texttt{Ovo je mojOPIS:} \\
\texttt{Mreža je dimenzija 5x5.} \\
\texttt{Gornji lijevi kut je prazan (▢), a donji desni kut je popunjen (X).} \\
\texttt{Popunjeni X{-}ovi tvore oblik obrnutih stepenica koje se spuštaju od gornjeg desnog kuta prema donjem lijevom kutu.} \\
\texttt{Redovi od vrha prema dnu:} \\
\texttt{1. red: 4 prazna, 1 X} \\
\texttt{2. red: 3 prazna, 2 X} \\
\texttt{3. red: 2 prazna, 3 X} \\
\texttt{4. red: 1 prazno, 4 X} \\
\texttt{5. red: 5 X} \\
\texttt{PITANJE:} \\
\texttt{U tvojoj mreži, je li donji lijevi kut (pozicija 5,1) prazan (▢) ili popunjen (X)?Započni svoj odgovor s "ODGOVOR:" i nemoj dodati ništa drugo.} \\
            }
        }
    }
    & \\ \\

    \theutterance \stepcounter{utterance}  
    & \multicolumn{4}{p{0.6\linewidth}}{
        \cellcolor[rgb]{0.9,0.9,0.9}{
            \makecell[{{p{\linewidth}}}]{
                \texttt{\tiny{[P1$\rangle$GM]}}
                \texttt{ODGOVOR: Popunjen (X).} \\
            }
        }
    }
    & & \\ \\

    \theutterance \stepcounter{utterance}  
    & & & \multicolumn{2}{p{0.3\linewidth}}{
        \cellcolor[rgb]{0.9,0.9,0.9}{
            \makecell[{{p{\linewidth}}}]{
                \texttt{\tiny{[GM$|$GM]}}
                \texttt{continue} \\
            }
        }
    }
    & & \\ \\

    \theutterance \stepcounter{utterance}  
    & & \multicolumn{4}{p{0.6\linewidth}}{
        \cellcolor[rgb]{0.9,0.9,0.9}{
            \makecell[{{p{\linewidth}}}]{
                \texttt{\tiny{[P1$\langle$GM]}}
                \texttt{Sada postavi pitanje kako bi saznao nove aspekte moje mreže koji se možda razlikuju od tvoje mreže. Započni s "PITANJE:" i ne dodavaj ništa drugo.} \\
            }
        }
    }
    & \\ \\

    \theutterance \stepcounter{utterance}  
    & \multicolumn{4}{p{0.6\linewidth}}{
        \cellcolor[rgb]{0.9,0.9,0.9}{
            \makecell[{{p{\linewidth}}}]{
                \texttt{\tiny{[P1$\rangle$GM]}}
                \texttt{PITANJE: Je li u tvojoj mreži središnje polje (pozicija 3,3) prazno (▢) ili popunjeno (X)?} \\
            }
        }
    }
    & & \\ \\

    \theutterance \stepcounter{utterance}  
    & & & \multicolumn{2}{p{0.3\linewidth}}{
        \cellcolor[rgb]{0.9,0.9,0.9}{
            \makecell[{{p{\linewidth}}}]{
                \texttt{\tiny{[GM$|$GM]}}
                \texttt{continue} \\
            }
        }
    }
    & & \\ \\

    \theutterance \stepcounter{utterance}  
    & & \multicolumn{4}{p{0.6\linewidth}}{
        \cellcolor[rgb]{0.9,0.9,0.9}{
            \makecell[{{p{\linewidth}}}]{
                \texttt{\tiny{[GM$\rangle$P2]}}
                \texttt{ODGOVOR: Popunjen (X).} \\
\texttt{PITANJE: Je li u tvojoj mreži središnje polje (pozicija 3,3) prazno (▢) ili popunjeno (X)?Započni svoj odgovor s "ODGOVOR:" i nemoj dodati ništa drugo.} \\
            }
        }
    }
    & \\ \\

    \theutterance \stepcounter{utterance}  
    & & & \multicolumn{4}{p{0.6\linewidth}}{
        \cellcolor[rgb]{0.9,0.9,0.9}{
            \makecell[{{p{\linewidth}}}]{
                \texttt{\tiny{[GM$\langle$P2]}}
                \texttt{ODGOVOR: Popunjeno (X).} \\
\texttt{PITANJE: U tvom drugom redu, jesu li polja (2,2) i (2,4) oba popunjena (X)?} \\
            }
        }
    }
     \\ \\

    \theutterance \stepcounter{utterance}  
    & & & \multicolumn{2}{p{0.3\linewidth}}{
        \cellcolor[rgb]{0.9,0.9,0.9}{
            \makecell[{{p{\linewidth}}}]{
                \texttt{\tiny{[GM$|$GM]}}
                \texttt{continue} \\
            }
        }
    }
    & & \\ \\

    \theutterance \stepcounter{utterance}  
    & & \multicolumn{4}{p{0.6\linewidth}}{
        \cellcolor[rgb]{0.9,0.9,0.9}{
            \makecell[{{p{\linewidth}}}]{
                \texttt{\tiny{[GM$\rangle$P2]}}
                \texttt{Sada postavi pitanje kako bi saznao nove aspekte moje mreže koji se možda razlikuju od tvoje mreže. Započni s "PITANJE:" i ne dodavaj ništa drugo.} \\
            }
        }
    }
    & \\ \\

    \theutterance \stepcounter{utterance}  
    & & & \multicolumn{4}{p{0.6\linewidth}}{
        \cellcolor[rgb]{0.9,0.9,0.9}{
            \makecell[{{p{\linewidth}}}]{
                \texttt{\tiny{[GM$\langle$P2]}}
                \texttt{PITANJE: U tvojoj mreži, u prvom redu, jesu li polja (1,2) i (1,4) ista (oba X ili oba prazna)?} \\
            }
        }
    }
     \\ \\

    \theutterance \stepcounter{utterance}  
    & & & \multicolumn{2}{p{0.3\linewidth}}{
        \cellcolor[rgb]{0.9,0.9,0.9}{
            \makecell[{{p{\linewidth}}}]{
                \texttt{\tiny{[GM$|$GM]}}
                \texttt{continue} \\
            }
        }
    }
    & & \\ \\

    \theutterance \stepcounter{utterance}  
    & & \multicolumn{4}{p{0.6\linewidth}}{
        \cellcolor[rgb]{0.9,0.9,0.9}{
            \makecell[{{p{\linewidth}}}]{
                \texttt{\tiny{[P1$\langle$GM]}}
                \texttt{ODGOVOR: Popunjeno (X).} \\
\texttt{PITANJE: U tvom drugom redu, jesu li polja (2,2) i (2,4) oba popunjena (X)?} \\
\texttt{PITANJE: U tvojoj mreži, u prvom redu, jesu li polja (1,2) i (1,4) ista (oba X ili oba prazna)?Započni svoj odgovor s "ODGOVOR:" i nemoj dodati ništa drugo.} \\
            }
        }
    }
    & \\ \\

    \theutterance \stepcounter{utterance}  
    & \multicolumn{4}{p{0.6\linewidth}}{
        \cellcolor[rgb]{0.9,0.9,0.9}{
            \makecell[{{p{\linewidth}}}]{
                \texttt{\tiny{[P1$\rangle$GM]}}
                \texttt{ODGOVOR: U drugom redu, polje (2,2) je prazno (▢), a polje (2,4) je prazno (▢). U prvom redu, polje (1,2) je X, a polje (1,4) je X (oba su X).} \\
            }
        }
    }
    & & \\ \\

    \theutterance \stepcounter{utterance}  
    & & & \multicolumn{2}{p{0.3\linewidth}}{
        \cellcolor[rgb]{0.9,0.9,0.9}{
            \makecell[{{p{\linewidth}}}]{
                \texttt{\tiny{[GM$|$GM]}}
                \texttt{continue} \\
            }
        }
    }
    & & \\ \\

    \theutterance \stepcounter{utterance}  
    & & \multicolumn{4}{p{0.6\linewidth}}{
        \cellcolor[rgb]{0.9,0.9,0.9}{
            \makecell[{{p{\linewidth}}}]{
                \texttt{\tiny{[P1$\langle$GM]}}
                \texttt{Sada postavi pitanje kako bi saznao nove aspekte moje mreže koji se možda razlikuju od tvoje mreže. Započni s "PITANJE:" i ne dodavaj ništa drugo.} \\
            }
        }
    }
    & \\ \\

    \theutterance \stepcounter{utterance}  
    & \multicolumn{4}{p{0.6\linewidth}}{
        \cellcolor[rgb]{0.9,0.9,0.9}{
            \makecell[{{p{\linewidth}}}]{
                \texttt{\tiny{[P1$\rangle$GM]}}
                \texttt{PITANJE: U tvojoj mreži, je li polje (2,5) (kraj drugog reda, desno) prazno (▢) ili popunjeno (X)?} \\
            }
        }
    }
    & & \\ \\

    \theutterance \stepcounter{utterance}  
    & & & \multicolumn{2}{p{0.3\linewidth}}{
        \cellcolor[rgb]{0.9,0.9,0.9}{
            \makecell[{{p{\linewidth}}}]{
                \texttt{\tiny{[GM$|$GM]}}
                \texttt{continue} \\
            }
        }
    }
    & & \\ \\

    \theutterance \stepcounter{utterance}  
    & & \multicolumn{4}{p{0.6\linewidth}}{
        \cellcolor[rgb]{0.9,0.9,0.9}{
            \makecell[{{p{\linewidth}}}]{
                \texttt{\tiny{[GM$\rangle$P2]}}
                \texttt{ODGOVOR: U drugom redu, polje (2,2) je prazno (▢), a polje (2,4) je prazno (▢). U prvom redu, polje (1,2) je X, a polje (1,4) je X (oba su X).} \\
\texttt{PITANJE: U tvojoj mreži, je li polje (2,5) (kraj drugog reda, desno) prazno (▢) ili popunjeno (X)?Započni svoj odgovor s "ODGOVOR:" i nemoj dodati ništa drugo.} \\
            }
        }
    }
    & \\ \\

    \theutterance \stepcounter{utterance}  
    & & & \multicolumn{4}{p{0.6\linewidth}}{
        \cellcolor[rgb]{0.9,0.9,0.9}{
            \makecell[{{p{\linewidth}}}]{
                \texttt{\tiny{[GM$\langle$P2]}}
                \texttt{ODGOVOR: Popunjeno (X).} \\
\texttt{PITANJE: U tvojoj mreži, u četvrtom redu, je li polje (4,2) prazno (▢)?} \\
            }
        }
    }
     \\ \\

    \theutterance \stepcounter{utterance}  
    & & & \multicolumn{2}{p{0.3\linewidth}}{
        \cellcolor[rgb]{0.9,0.9,0.9}{
            \makecell[{{p{\linewidth}}}]{
                \texttt{\tiny{[GM$|$GM]}}
                \texttt{continue} \\
            }
        }
    }
    & & \\ \\

    \theutterance \stepcounter{utterance}  
    & & \multicolumn{4}{p{0.6\linewidth}}{
        \cellcolor[rgb]{0.9,0.9,0.9}{
            \makecell[{{p{\linewidth}}}]{
                \texttt{\tiny{[GM$\rangle$P2]}}
                \texttt{Sada postavi pitanje kako bi saznao nove aspekte moje mreže koji se možda razlikuju od tvoje mreže. Započni s "PITANJE:" i ne dodavaj ništa drugo.} \\
            }
        }
    }
    & \\ \\

    \theutterance \stepcounter{utterance}  
    & & & \multicolumn{4}{p{0.6\linewidth}}{
        \cellcolor[rgb]{0.9,0.9,0.9}{
            \makecell[{{p{\linewidth}}}]{
                \texttt{\tiny{[GM$\langle$P2]}}
                \texttt{PITANJE: U tvojoj mreži, u petom redu, jesu li sva polja popunjena (X)?} \\
            }
        }
    }
     \\ \\

    \theutterance \stepcounter{utterance}  
    & & & \multicolumn{2}{p{0.3\linewidth}}{
        \cellcolor[rgb]{0.9,0.9,0.9}{
            \makecell[{{p{\linewidth}}}]{
                \texttt{\tiny{[GM$|$GM]}}
                \texttt{continue} \\
            }
        }
    }
    & & \\ \\

    \theutterance \stepcounter{utterance}  
    & & \multicolumn{4}{p{0.6\linewidth}}{
        \cellcolor[rgb]{0.9,0.9,0.9}{
            \makecell[{{p{\linewidth}}}]{
                \texttt{\tiny{[P1$\langle$GM]}}
                \texttt{ODGOVOR: Popunjeno (X).} \\
\texttt{PITANJE: U tvojoj mreži, u četvrtom redu, je li polje (4,2) prazno (▢)?} \\
\texttt{PITANJE: U tvojoj mreži, u petom redu, jesu li sva polja popunjena (X)?Započni svoj odgovor s "ODGOVOR:" i nemoj dodati ništa drugo.} \\
            }
        }
    }
    & \\ \\

    \theutterance \stepcounter{utterance}  
    & \multicolumn{4}{p{0.6\linewidth}}{
        \cellcolor[rgb]{0.9,0.9,0.9}{
            \makecell[{{p{\linewidth}}}]{
                \texttt{\tiny{[P1$\rangle$GM]}}
                \texttt{ODGOVOR: U četvrtom redu, polje (4,2) je popunjeno (X). U petom redu, sva polja su popunjena (X).} \\
            }
        }
    }
    & & \\ \\

    \theutterance \stepcounter{utterance}  
    & & & \multicolumn{2}{p{0.3\linewidth}}{
        \cellcolor[rgb]{0.9,0.9,0.9}{
            \makecell[{{p{\linewidth}}}]{
                \texttt{\tiny{[GM$|$GM]}}
                \texttt{continue} \\
            }
        }
    }
    & & \\ \\

    \theutterance \stepcounter{utterance}  
    & & \multicolumn{4}{p{0.6\linewidth}}{
        \cellcolor[rgb]{0.9,0.9,0.9}{
            \makecell[{{p{\linewidth}}}]{
                \texttt{\tiny{[P1$\langle$GM]}}
                \texttt{Sada postavi pitanje kako bi saznao nove aspekte moje mreže koji se možda razlikuju od tvoje mreže. Započni s "PITANJE:" i ne dodavaj ništa drugo.} \\
            }
        }
    }
    & \\ \\

    \theutterance \stepcounter{utterance}  
    & \multicolumn{4}{p{0.6\linewidth}}{
        \cellcolor[rgb]{0.9,0.9,0.9}{
            \makecell[{{p{\linewidth}}}]{
                \texttt{\tiny{[P1$\rangle$GM]}}
                \texttt{PITANJE: U tvojoj mreži, je li polje (1,1) (gornji lijevi kut) prazno (▢) ili popunjeno (X)?} \\
            }
        }
    }
    & & \\ \\

    \theutterance \stepcounter{utterance}  
    & & & \multicolumn{2}{p{0.3\linewidth}}{
        \cellcolor[rgb]{0.9,0.9,0.9}{
            \makecell[{{p{\linewidth}}}]{
                \texttt{\tiny{[GM$|$GM]}}
                \texttt{continue} \\
            }
        }
    }
    & & \\ \\

    \theutterance \stepcounter{utterance}  
    & & \multicolumn{4}{p{0.6\linewidth}}{
        \cellcolor[rgb]{0.9,0.9,0.9}{
            \makecell[{{p{\linewidth}}}]{
                \texttt{\tiny{[GM$\rangle$P2]}}
                \texttt{ODGOVOR: U četvrtom redu, polje (4,2) je popunjeno (X). U petom redu, sva polja su popunjena (X).} \\
\texttt{PITANJE: U tvojoj mreži, je li polje (1,1) (gornji lijevi kut) prazno (▢) ili popunjeno (X)?Započni svoj odgovor s "ODGOVOR:" i nemoj dodati ništa drugo.} \\
            }
        }
    }
    & \\ \\

    \theutterance \stepcounter{utterance}  
    & & & \multicolumn{4}{p{0.6\linewidth}}{
        \cellcolor[rgb]{0.9,0.9,0.9}{
            \makecell[{{p{\linewidth}}}]{
                \texttt{\tiny{[GM$\langle$P2]}}
                \texttt{ODGOVOR: Prazno (▢).} \\
\texttt{PITANJE: U tvojoj mreži, u drugom redu, jesu li polja (2,1) i (2,3) oba popunjena (X)?} \\
            }
        }
    }
     \\ \\

    \theutterance \stepcounter{utterance}  
    & & & \multicolumn{2}{p{0.3\linewidth}}{
        \cellcolor[rgb]{0.9,0.9,0.9}{
            \makecell[{{p{\linewidth}}}]{
                \texttt{\tiny{[GM$|$GM]}}
                \texttt{continue} \\
            }
        }
    }
    & & \\ \\

    \theutterance \stepcounter{utterance}  
    & & \multicolumn{4}{p{0.6\linewidth}}{
        \cellcolor[rgb]{0.9,0.9,0.9}{
            \makecell[{{p{\linewidth}}}]{
                \texttt{\tiny{[GM$\rangle$P2]}}
                \texttt{Sada donesi odluku. Što misliš: jesu li tvoja mreža i druga opisana mreža iste? Napiši "ODLUKA: ista mreža" ako misliš da su iste mreže ili "ODLUKA: različite mreže" ako misliš da su različite mreže. Ne dodaj ništa drugo.} \\
            }
        }
    }
    & \\ \\

    \theutterance \stepcounter{utterance}  
    & & & \multicolumn{4}{p{0.6\linewidth}}{
        \cellcolor[rgb]{0.9,0.9,0.9}{
            \makecell[{{p{\linewidth}}}]{
                \texttt{\tiny{[GM$\langle$P2]}}
                \texttt{ODLUKA: različite mreže} \\
            }
        }
    }
     \\ \\

    \theutterance \stepcounter{utterance}  
    & & & \multicolumn{2}{p{0.3\linewidth}}{
        \cellcolor[rgb]{0.9,0.9,0.9}{
            \makecell[{{p{\linewidth}}}]{
                \texttt{\tiny{[GM$|$GM]}}
                \texttt{continue} \\
            }
        }
    }
    & & \\ \\

    \theutterance \stepcounter{utterance}  
    & & & \multicolumn{2}{p{0.3\linewidth}}{
        \cellcolor[rgb]{0.9,0.9,0.9}{
            \makecell[{{p{\linewidth}}}]{
                \texttt{\tiny{[GM$|$GM]}}
                \texttt{success} \\
            }
        }
    }
    & & \\ \\

    \theutterance \stepcounter{utterance}  
    & & \multicolumn{4}{p{0.6\linewidth}}{
        \cellcolor[rgb]{0.9,0.9,0.9}{
            \makecell[{{p{\linewidth}}}]{
                \texttt{\tiny{[P1$\langle$GM]}}
                \texttt{ODGOVOR: Prazno (▢).} \\
\texttt{PITANJE: U tvojoj mreži, u drugom redu, jesu li polja (2,1) i (2,3) oba popunjena (X)?} \\
\texttt{Sada donesi odluku. Što misliš: jesu li tvoja mreža i druga opisana mreža iste? Napiši "ODLUKA: ista mreža" ako misliš da su iste mreže ili "ODLUKA: različite mreže" ako misliš da su različite mreže. Ne dodaj ništa drugo.} \\
            }
        }
    }
    & \\ \\

    \theutterance \stepcounter{utterance}  
    & \multicolumn{4}{p{0.6\linewidth}}{
        \cellcolor[rgb]{0.9,0.9,0.9}{
            \makecell[{{p{\linewidth}}}]{
                \texttt{\tiny{[P1$\rangle$GM]}}
                \texttt{ODGOVOR: U drugom redu, polje (2,1) je popunjeno (X), a polje (2,3) je popunjeno (X).} \\
\texttt{ODLUKA: različite mreže} \\
            }
        }
    }
    & & \\ \\

    \theutterance \stepcounter{utterance}  
    & & & \multicolumn{2}{p{0.3\linewidth}}{
        \cellcolor[rgb]{0.9,0.9,0.9}{
            \makecell[{{p{\linewidth}}}]{
                \texttt{\tiny{[GM$|$GM]}}
                \texttt{abort, first line: ODGOVOR: U drugom redu, polje (2,1) je popunjeno (X), a polje (2,3) je popunjeno (X)., but expected ODLUKA:} \\
            }
        }
    }
    & & \\ \\

\end{supertabular}
}

\end{document}
